\documentclass[landscape,a1paper]{baposter}

\usepackage[portuguese]{babel}
\usepackage[utf8]{inputenc}
\usepackage{graphicx}
\usepackage{helvet}
\renewcommand{\familydefault}{\sfdefault}

% Necessário para que o BibTeX funcione com o Baposter
% Esse arquivo precisa estar na mesma pasta do poster.tex
% Baixe em:
% http://www.tex.ac.uk/tex-archive/biblio/bibtex/contrib/authordate/authordate1-4.sty
\usepackage{authordate1-4}

% The 'caption' package provides a nicer-looking replacement
\usepackage[labelfont=bf,textfont=it]{caption}

% Habilita subfigure.
\usepackage{subcaption}

% Habilita o uso de matrix.
\usepackage{amsmath}

% Habilita ambiente de algoritmos.
\usepackage[linesnumbered,portuguese]{algorithm2e}

% Divide igualmente o conteúdo da última página entre as duas colunas.
%\usepackage{flushend}

\usepackage{url}

% Copiado do template original:
\usepackage{relsize}
\usepackage{multirow}
\usepackage{booktabs}
\usepackage{multicol}
\usepackage{ae}

% Permite colocar direto o nome da imagem no includegraphics.
\graphicspath{{images/}}

% Distância entre colunas.
\setlength{\columnsep}{0.7em}

% Espessura da linha vertical que separa as colunas.
\setlength{\columnseprule}{0mm}

%
% Challenges and Goals
% =====================
%
% Education
% ------
%
% - Computer literacy is more than just using computers, but also programming them
%     > similar metaphor to non-professional writers learning how to write
% - The skills involved in computer programming are useful for many other activitites
%     > thinking about thinking
%     > trial and error
% - Existing solutions are too focused on computer languages or too far from the
%   Brazilian culture
%     > brazilian state division and fauna
%
% Game
% ------
%
% - A puzzle in which players need to experimentaly design an ice cream factory to
%   produce the requested flavors
% - Simple mechanics based on path decisions
%     > with switchers and scales to conduct the products on the conveyors
% - Fantasy and aesthetics with a cartoonished style and funny animations
%
% Prototype
% ------
%
% - Experimentation of the basic mechanics and aesthetics, and interaction methods
% - Presentation of the visual identity and main characters
%     > penguin (that comes from Argentina to the south region of Brazil) and
%       jaguar (that lives in the northest region of Brazil)
%
% Research
% ------
%
% - To be used for complementing the computer programming classes in the
%   Administration course from the University of São Paulo
% - To be used as a sample game in a PhD research interested in evaluating the
%   educational potential of games **** (to falando sobre a pesquisa do Bosco.
%   Talvez fosse legal perguntar pra ele - e pro professor! - sobre como escrever
%   isso aqui)
%
% Fun
% ------
%
% The game aims to produce a fun experience through different aspects:
% - Fantasy: an ice cream factory managed by the player with the help of penguin
%   assistants and other animals as truck drivers
% - Challenge: different configurations of conveyors in many levels, producing
%   initially simple problems but increasing complexity as the player progress in
%   the game
% - "Juiceness": pleasing designs and animations, with cartoonish and comical
%   styles strongly based on visual and auditory puns and incongruent elements
%   ("crazy machines")
% - Curiosity: elicitation of curiosity regarding the appearance and behavior of
%   new animals and devices (truck drivers and control devices unlocked as the
%   player progress in the game)
%
% Target
% ------
%
% - Children with at least 12 years old and adults
% - Playable in Windows, Linux, Android and iOS platforms
%     > and also via Web through Flash or HTML5
%
% Game Mechanics and Aesthetics
% =====================
%
% **** A ideia de construir um diagrama aqui com o "processo de jogo" é bastante
% bacana. Acho que aqui dá de colocar também algumas imagens do Danilo.
%
% Production
% =====================
%
% Engine
% ------
%
% - Currently using the HaxeFlixel engine
%     > open source engine based on the ActionScript 3 programming language
%       (Haxe + Flixel products)
%     > able to compile to all targeted platforms
%

\begin{document}
\begin{poster}{
    grid=false,
    eyecatcher=true, % Exibe imagem à esquerda do título.
    colspacing=0.7em,
    %%% Definindo as cores das caixas. %%%
    headerColorOne=black!45!cyan!15,
    borderColor=black!45!cyan!15,
    %%% Estilo da borda. %%%
    textborder=faded,
    %%% Estilo da sub-caixa de título de cada caixa. %%%
    headerborder=open,
    headershape=roundedright,
    headershade=plain,
    background=none,
    bgColorOne=cyan!10!white,
    boxheaderheight=2.35em,
    %%% Espaço reservado para o título, autores e imagens superiores. %%%
    headerheight=11em,
    %%% Número de colunas (de 1 a 6). %%%
    columns=5
}
% Imagem superior esquerda.
{
    \includegraphics[height=10em]{sweetswitches}
}
% Título do pôster.
{
    \hfill
    \begin{minipage}{0.97\textwidth}
    In Sweet Switches players will face a series of puzzles inside ice cream
    factories. The objective is to attach devices to conveyor belts in order to
    respond to a production request. As those devices are analogous to
    programming statements, the game foster computational thinking.
    \end{minipage}
}
% Autores e nome da universidade.
{
    \scriptsize
    %\hfill%
    \centering
    \begin{minipage}{0.7\textwidth}
        \vskip 1.5em%
        %\centering
        \begin{minipage}{0.8\textwidth}
            Vinícius K. Daros, Luiz C. Vieira, Adalberto B. C. Pereira, Flávio S. C. da Silva\\[0.5em]
            Laboratory of Interactivity and Digital Entertainment Technology (LIDET)
            \\[0.15em]
            University of São Paulo - Brazil\\[0.5em]
            \texttt{\{vkdaros, lvieira, bosco, fcs\}@ime.usp.br}
        \end{minipage}
        \begin{minipage}{5em}
            \parbox[t][5em][t]{5em}{
                \begin{tabular}{l}
                    \includegraphics[width=5em]{usp-logo}\\
                    \includegraphics[width=5em]{universidade}
                \end{tabular}
            }
        \end{minipage}
    \end{minipage}

}
% Imagem superior direita.
{
%    \parbox[t][7.2em][b]{4.5em}{
%        \includegraphics[height=5.75em]{logo-ime}
%    }
}

\headerbox{Resumo}{name=abstract, column=0, row=0, span=1}{
    Conhecer a forma da pista pode dar muitos benefícios à Inteligência
    Artificial que controla os carros nos jogos de corrida, mas esse tipo de
    informação nem sempre está disponível ao NPC. Tendo como plataforma de
    testes o simulador de corridas TORCS, mostramos neste trabalho um método
    capaz de construir um modelo de pista baseado apenas em leituras de um
    número limitado de sensores no carro de corrida. Também apresentamos uma
    técnica de segmentação do circuito indicando o início e fim de trechos
    distintos, classificando-os por suas características - como reta, curva
    suave, curva fechada, entre outros.
}

\headerbox{End}{name=end, column=4, row=0, span=1}{
    Conhecer a forma da pista pode dar muitos benefícios à Inteligência
    Artificial que controla os carros nos jogos de corrida, mas esse tipo de
    informação nem sempre está disponível ao NPC. Tendo como plataforma de
    testes o simulador de corridas TORCS, mostramos neste trabalho um método

    \begin{minipage}{\textwidth}
        \centering
        \vskip 0.25em
        \includegraphics[width=0.9\textwidth]{onca}
    \end{minipage}
}

\headerbox{Links}{name=links, column=3, below=end, span=2}{
    \url{http://www.ime.usp.br/~lidet/sweet-switches/}
    \url{http://github.com/vkdaros/igamer2014}
}

\headerbox{Construindo}{name=model, column=1, aligned=abstract, span=3}{
    \begin{multicols}{3}
        Durante uma volta de reconhecimento, o piloto virtual coleta dados
        através de sensores, como se fosse um robô. Ele tenta se manter no
        centro da pista, anda relativamente devagar (por volta de $60 Km/h$ em
        retas e $20 Km/h$ em curvas) e armazena a leitura dos sensores a cada 5
        metros. Ao fim da volta, os dados são processados. Assim, considerando
        pares consecutivos de medições, temos pequenos \textbf{segmentos} da
    \end{multicols}
}

\headerbox{Construindo2}{name=bot, column=1, above=bottom, span=2}{
    \begin{multicols}{3}
        Durante uma volta de reconhecimento, o piloto virtual coleta dados
        através de sensores, como se fosse um robô. Ele tenta se manter no
        centro da pista, anda relativamente devagar (por volta de $60 Km/h$ em
        retas e $20 Km/h$ em curvas) e armazena a leitura dos sensores a cada 5
        metros. Ao fim da volta, os dados são processados. Assim, considerando
        pares consecutivos de medições, temos pequenos \textbf{segmentos} da
    \end{multicols}
}

\end{poster}
\end{document}
